\section{推测执行留下的微架构痕迹}

错误路径上的指令不会退休，因此不会把寄存器结果或 store 变成架构可见状态；但是“没有退休”并不意味着“没有发生过任何硬件动作”。推测 load 可能填充 Cache line，分支可能更新预测器状态，页表遍历可能占用 TLB 或 walk cache，执行端口也可能形成可测量的争用。攻击者若能反复控制输入并测量这些状态差异，就可能从时间相关性中推断本来无权读取的数据。Spectre 类问题利用的正是“架构状态已回滚、微架构痕迹仍可观察”这条边界。\cite{linux-speculation,linux-spectre}

\topic{一条边界检查的瞬态执行 trace}
设程序先检查索引 $i<N$，再读取数组元素。训练阶段使预测器倾向于“检查通过”；攻击阶段给出越界索引，前端可能在比较结果返回前沿预测路径发出 load。权限检查和分支解析最终阻止错误结果退休，但该 load 及其后续依赖访问可能已经改变 Cache 状态。攻击者随后测量一组合法地址的访问时间，并通过命中与未命中的差异推断瞬态路径使用过哪个值。

这个过程有六个必须分开的事件：

\begin{enumerate}[leftmargin=2.2em]
  \item 预测器选择路径，它只决定暂时取哪些指令；
  \item load 形成地址并请求存储层次；
  \item 权限、依赖和分支条件在各自时刻得到解析；
  \item 失败检查阻止结果退休，并触发流水线恢复；
  \item Cache 或预测器中的部分状态没有随 ROB 一起恢复；
  \item 测量代码把时延差转换为可观察信息。
\end{enumerate}

因此，“处理器支持精确异常”只能证明架构状态正确，不能单独证明不存在信息泄漏。

\topic{缓解措施必须对应具体的预测与共享状态}
在不可信索引进入敏感访存前放置体系结构规定的推测屏障，可以阻止后续访问越过检查；对间接分支，可使用受控预测、预测器清理或编译器生成的受限跳转序列；在安全域切换处，还可能需要清理或隔离分支历史。常量时间算法避免让秘密决定访存地址或控制流，它处理的是信息编码方式，不等同于清空预测器。Cache 分区、核心调度和限制 SMT 则处理共享资源边界，也不能替代正确的边界检查。

任何缓解措施都不应只写成一个名称。验证时要明确：它约束哪一类推测，在哪个执行点生效，是否跨逻辑线程或特权级，能否覆盖当前处理器型号，以及会增加多少屏障、预测失败、Cache miss 或上下文切换开销。

\topic{性能代价要在威胁边界内衡量}
完全禁止预测会显著降低现代流水线利用率，因此系统通常按信任域和工作负载选择措施。内核与用户态切换、虚拟机切换、浏览器沙箱和密码学例程面对的攻击面不同。正确做法是先标出秘密、攻击者可控输入、可共享状态和可用计时源，再选择最窄但足够的隔离措施。若只比较平均运行时间而不验证安全域之间是否仍可形成相关时延，性能测试无法证明缓解有效。

\begin{warningbox}
将一次未退休的 load 称为“没有执行”会混淆执行与提交。它没有产生架构可见结果，却可能已经请求 Cache、占用 MSHR 或影响预测状态。分析安全问题时必须同时画出 ROB 提交边界和微架构状态边界。
\end{warningbox}
